\section{Trace and Norm}

\begin{definition}[\textbf{Trace and Norm}]
	Fix a field \( K \), let \( A \) being a commutative \( K \)-algebra, with \( K \to A \), s.t. \(
	\dim_K(A) < \infty \). Given \( a \in A \), let
	\[
		\begin{aligned}
			\vp_a : A & \to A \\
			\vp_a(b)  & = ab
		\end{aligned}
	\]
	this is a \( K \)-linear map.
	\begin{itemize}
		\item The \underline{trace} \( \Tr_{\qo{A}{K}}(a) \) is trace(\( \vp_a \)) \( \in K \).
		\item The \underline{norm} \( N_{\qo{A}{K}}(a) \) is \( \det(\vp_a) \in K \).
	\end{itemize}
\end{definition}

\begin{eg}
	\leavevmode
	\begin{enumerate}
		\item If \( a \in K \implies \) w.r.t. any basis of \( A \), \( \vp_a \) correspond to \( a I_n
		      \implies \)
		      \[
			      \begin{aligned}
				      \Tr_{\qo{A}{K}}(a) & = na  \\
				      N_{\qo{A}{K}}(a)   & = a^n
			      \end{aligned}
		      \]
		      with \( n = \dim_{K}(A) \).
		\item Let \( QQ \hookrightarrow L = \QQ(\sqrt{d}) \) with \( d \in \ZZ \) not a square. We
		      have basis \( \qo{L}{\QQ} \) being \( 1, \sqrt{d} \). Now if \( a = a_1 + a_2 \sqrt{d} \),
		      we have:
		      \[
			      \begin{aligned}
				      \begin{cases}
					      a\cdot 1 = a_1 + a_2 \sqrt{d} \\
					      a\cdot \sqrt{d} = d a_2 + a_1 \sqrt{d}
				      \end{cases}
				       & \implies
				      \begin{pmatrix}
					      a_1  & a_2 \\
					      da_2 & a_1
				      \end{pmatrix} \\
				       & \implies
				      \begin{cases}
					      \Tr_{\qo{L}{\QQ}}(a) = 2a_1 \\
					      N_{\qo{L}{\QQ}}(a) = a_1^2 -da_2^2 = (a_1 + \sqrt{d}a_2) (a_1 - \sqrt{d}a_2)
				      \end{cases}
			      \end{aligned}
		      \]
	\end{enumerate}
\end{eg}

\begin{proposition}[\textbf{Easy Properties}]
	\leavevmode
	\begin{enumerate}
		\item \( \Tr_{\qo{A}{K}}: A \to K \) is a \( K \)-linear map.
		      \begin{proof}
			      Have
			      \[
				      \vp_{\lambda_1 a_1 + \lambda _2 a_2} = \lambda _1 \vp_{a_1} + \lambda _2 \vp_{a_2} \quad
				      a_1,a_2 \in A, \lambda _1, \lambda _2 \in K
			      \]
			      and trace(-) is a \( K \)-linear map.
		      \end{proof}
		\item \( N_{\qo{A}{K}}(ab) = N_{\qo{A}{K}}(a) N_{\qo{A}{K}}(b) \; \forall \; a,b \in A \).
		      \begin{proof}
			      See that
			      \[
				      \vp_{ab} = \vp_a \circ \vp_b
			      \]
			      and \( \det(-) \) is multiplicative.
		      \end{proof}
		\item (Transitivity) \( \underbrace{k\hookrightarrow }_{\text{finite ext.}} K \to A \)
		      \begin{itemize}
			      \item \( \Tr_{\qo{A}{K}}(a) = \Tr_{\qo{K}{k}}(\Tr_{\qo{A}{K}} (a)) \).
			      \item \( N_{\qo{A}{k}}(a) = N_{\qo{K}{k}}(N_{\qo{A}{K}}(a)) \).
		      \end{itemize}
		      We only prove for the first, the second is similar.
		      \begin{proof}
			      Say
			      \[
				      \begin{aligned}
					      e_1, & \ldots, e_n \text{ basis of } \qo{A}{K} \\
					      f_1, & \ldots, f_m \text{ basis of } \qo{K}{k}
				      \end{aligned}
			      \]
			      we have a basis given by
			      \[
				      \{e_i f_j \; | \; i \leq n, j \leq m\}
			      \]
			      for \( \qo{A}{k} \).\todo{chech this.} Now given \( a\in A \), have:
			      \[
				      \begin{aligned}
					      a e_i                                         & = \sum_{j=1}^n \alpha_{ji} e_j \quad \alpha_{ji} \in K \\
					      \alpha_{ji} f_p                               & = \sum_{q=1}^m \underbrace{\gamma_{jiqp}}_{\in k} f_q  \\
					      \implies a e_i f_p                            & = \sum_{\substack{j = 1 \sim n                         \\ q = 1\sim m}}
					      \gamma_{jiqp} e_j f_p \implies \Tr_{\qo{A}{k}}(a) = \sum_{\substack{j = 1 \sim n                       \\
					      q = 1 \sim m}} \gamma_{jiqp}                                                                           \\
					      \Tr_{\qo{A}{K}}(a)                            & = \sum_{j=1}^n \alpha_{jj}                             \\
					      \implies \Tr_{\qo{K}{k}}(\Tr_{\qo{A}{K}} (a)) & = \sum_{j=1}^n
					      \Tr_{\qo{K}{k}}(\alpha_{jj})                                                                           \\
					                                                    & = \sum_{j=1}^n \sum_{q=1}^m \gamma_{jjqq}
				      \end{aligned}
			      \]
			      which completes the proof.
		      \end{proof}
		\item Suppose the following diagram:
		      \[
			      \begin{tikzcd}
				      K \arrow[r] \arrow[d, hook] & A \arrow[d] \\
				      L \arrow[r] & A_L = A \otimes_K L
			      \end{tikzcd}
		      \]
		      let \( a\in A \implies \)
		      \begin{itemize}
			      \item \( \Tr_{\qo{A}{K}}(a) = \Tr_{\qo{A_L}{L}}(a\otimes 1) \).
			      \item \( N_{\qo{A}{K}}(a) = N_{\qo{A_L}{L}}(a\otimes 1) \).
		      \end{itemize}

		      \begin{proof}
			      If \( e_1,\ldots, e_n \) be basis of \( \qo{A}{K} \), then \( e_1\otimes 1,\ldots, e_n
			      \otimes 1 \) is basis of \( \qo{A_L}{L} \) \todo{check this!}. Notice that
			      \[
				      (a\otimes 1)(e_i \otimes 1) = ae_i \otimes 1
			      \]
			      multiply by \( a \), resp \( a\otimes 1 \) is given by the same matrix.
		      \end{proof}
		\item If \( A = A_1 \times A_2 \), then
		      \begin{itemize}
			      \item \( \Tr_{\qo{A}{K}}(a_1,a_2) = \Tr_{\qo{A_1}{K}}(a_1) + \Tr_{\qo{A_2}{K}}(a_2) \).
			      \item \( N_{\qo{A}{K}}(a_1,a_2) = N_{\qo{A_1}{K}}(a_1) \cdot N_{\qo{A_2}{K}}(a_2) \).
		      \end{itemize}

		      \begin{proof}
			      Given basis for \( A_1,A_2 \rsa \) basis for \( A_1\times A_2 \). If \( M_1, M_2 \) are
			      matrices corresponds to multiplication by \( a_1,a_2 \) in \( A_1,A_2 \) respectively,
			      then the matrix of multiplication by \( (a_1,a_2) \) is given by
			      \[
				      \begin{pmatrix}
					      M_1 & 0   \\
					      0   & M_2
				      \end{pmatrix}
			      \]
		      \end{proof}
	\end{enumerate}
\end{proposition}
Now we want to discuss how the notion of trace and norm related to separable extension?

\begin{lemma}
	\label{lem:1}
	If \( K \hookrightarrow L \) being field extension, \( a \in L \) is algebraic and separable over
	\( K \), let \( K' = K(a) \), then
	\[
		\Tr_{\qo{K'}{K}}: K' \to K \text{ is nonzero, hence surjective.}
	\]
\end{lemma}

\begin{proof}
	Let \( \overline{K} \) be the algebraic closure of \( K \), it is enough to show that
	\[
		\Tr_{\qo{A_{\overline K}}{\overline K}} \ne 0
	\]
	have
	\[
		\begin{tikzcd}
			K \arrow[r, hook] \arrow[d, hook] & K' = A \arrow[d] \\
			\overline{K} \arrow[r, hook] & A_{\overline{K}} = \overline{K} \otimes_K A
		\end{tikzcd}
	\]
	This follows since by properties 1 and 4 above:
	\[
		\Tr_{\qo{A_{\overline K}}{\overline K}} \bace{\sum_{i=1}^m a_i \otimes b_i} = \sum_{i=1}^n a_i
		\Tr_{\qo{K'}{K}}(b_i)
	\]
	suppose \( f \) is the minimal polynomial of \( a \) over \( K \), then \todo{check this!}
	\[
		\begin{aligned}
			K'              & \cong \qo{K[x]}{(f)}                                        \\
			A_{\overline K} & = \overline{K} \otimes_K K' \cong \qo{\overline{K}[x]}{(f)}
		\end{aligned}
	\]
	In \( \overline{K}[x] \), we have \( f = (x-\lambda_1)\cdots (x-\lambda_n) \), since \( a \) is
	separable, then \( \lambda_i \ne \lambda_j \) if \( i \ne j \), thus
	\[
		\qo{\overline{K}[x]}{(x-\lambda_1)\cdots(x-\lambda_n)} \cong \prod_{i=1}^n \frac{\overline{K}[x]}{(x-\lambda_i)}
		\cong \overline K \times \cdots \times \overline K
	\]
	By property 5, we have
	\[
		\Tr_{\qo{A_{\overline{K}}}{\overline K}} (a_1,\ldots, a_n) = a_1 + \cdots a_n \ne 0
	\]
\end{proof}

\begin{lemma}
	\label{lem:2}
	Let \( \text{char}(K) = p > 0 \). If \( K \hookrightarrow K' = K(a) \), where \( a^p \in K, a\not\in K \), then
	\[
		\Tr_{\qo{K'}{K}}: K' \to K \text{ is the 0 map.}
	\]
\end{lemma}

